%%% Please fill in basic information on your thesis, which will be automatically
%%% inserted at the right places. You need to replace \xxx{...} by real data.

% Type of your thesis:
%	"bc" for Bachelor's
%	"mgr" for Master's
%	"phd" for PhD
%	"rig" for rigorosum
\def\ThesisType{bc}

% Language of your study programme:
%	"cs" for Czech
%	"en" for English
\def\StudyLanguage{cs}

% Thesis title in English (exactly as in the official assignment)
% (Note: \xxx is a "ToDo label" which makes the unfilled visible. Remove it.)
\def\ThesisTitle{Arduino platform simulator}

% Author of the thesis (you)
\def\ThesisAuthor{Adam Balko}

% Year when the thesis is submitted
\def\YearSubmitted{2026}

% Name of the department or institute, where the work was officially assigned
% (according to the Organizational Structure of MFF UK in English,
% see https://www.mff.cuni.cz/en/faculty/organizational-structure,
% or a full name of a department outside MFF)
\def\Department{Department of Software and Computer Science Education}

% Is it a department (katedra), or an institute (ústav)?
\def\DeptType{Department}

% Thesis supervisor: name, surname and titles
\def\Supervisor{RNDr. David Bednárek, Ph.D.}

% Supervisor's department (again according to Organizational structure of MFF)
\def\SupervisorsDepartment{Department of Software Engineering}

% Study programme (does not apply to rigorosum theses)
\def\StudyProgramme{Computer Graphics, Vision and Computer Games Development}

% An optional dedication: you can thank whomever you wish (your supervisor,
% consultant, who provided you with tea and pizza, etc.)
\def\Dedication{%
\xxx{Dedication.}
}

% Abstract (recommended length around 80-200 words; this is not a copy of your thesis assignment!)
\def\Abstract{%
This thesis focuses simulating code of Arduino microcomputers in a virtual environment with a graphical user interface without the need of emulation. An important part of the thesis shows how and why should be the simulator and the GUI split into two processes, and how to keep them synchronised. This introduces the reader into the field of interprocess communication and its implementations using TCP connection, shared memory and events. Another important part focuses on simulation of digital circuits with extensible component library. It describes how new components and circuits can be created from YAML, SVG and C\# files by the user.
}

% 3 to 5 keywords (recommended) separated by \sep
% Keywords are useful for indexing and searching for the theses by topic.
\def\ThesisKeywords{%
digital circuit simulator, Arduino
}

% If any of your metadata strings contains TeX macros, you need to provide
% a plain-text version for use in XMP metadata embedded in the output PDF file.
% If you are not sure, check the generated thesis.xmpdata file.
\def\ThesisAuthorXMP{\ThesisAuthor}
\def\ThesisTitleXMP{\ThesisTitle}
\def\ThesisKeywordsXMP{\ThesisKeywords}
\def\AbstractXMP{\Abstract}

% If your abstracts are long and do not fit in the infopage, you can make the
% fonts a bit smaller by this setting. (Also, you should try to compress your abstract more.)
\def\InfoPageFont{}
%\def\InfoPageFont{\small}  % uncomment to decrease font size

% If you are studing in a Czech programme, you also need to provide metadata in Czech:
% (in English programmes, this is not used anywhere)

\def\ThesisTitleCS{Simulátor prostředí Arduino}
\def\DepartmentCS{Katedra softwaru a výuky informatiky}
\def\DeptTypeCS{Katedra}
\def\SupervisorsDepartmentCS{Katedra softwarového inženýrství}
\def\StudyProgrammeCS{Počítačová grafika, vidění a vývoj her}

\def\ThesisKeywordsCS{%
simulátor číslicových obvodů, Arduino
}

\def\AbstractCS{%
Tato práce se zaměřuje na simulaci kódu mikropočítačů Arduino ve virtuálním prostředí s grafickým uživatelským rozhraním bez nutnosti emulace. Důležitá část práce ukazuje, jak a proč by měly být simulátor a GUI rozděleny do dvou procesů a jak je udržovat synchronizované. To čtenáře uvádí do oblasti meziprocesové komunikace a její implementace pomocí TCP spojení, sdílené paměti a událostí. Další významná část se zaměřuje na simulaci digitálních obvodů s rozšiřitelnou knihovnou komponent. Popisuje, jak může uživatel vytvářet nové komponenty a obvody ze souborů YAML, SVG a C\#.
}
